\documentclass[11pt]{article}
\usepackage[a4paper,margin=25mm]{geometry}
\usepackage[T1]{fontenc}
\usepackage{lmodern,amsmath,amssymb,amsthm,array,booktabs}
\usepackage[hidelinks]{hyperref}
\newtheorem{theorem}{Theorem}
\newtheorem{lemma}[theorem]{Lemma}
\newtheorem{corollary}[theorem]{Corollary}
\newcommand{\F}{\mathbb F_2}
\newcommand{\rk}{\operatorname{rank}}
\newcommand{\Row}{\operatorname{Row}}
\newcommand{\GL}{\operatorname{GL}}
\setlength{\parskip}{4pt}
\setlength{\parindent}{0pt}
\title{Khesin--Li (2024), Section 8, Conjecture 8.3:\\
An elementary proof and a sharp correction to Proposition 8.2}
\author{Haobo Ma\\\small ChronoAI Pte Ltd, Singapore
\and Wenlin Zhang\\\small National University of Singapore\\\small ChronoAI Pte Ltd, Singapore}
\date{6 November 2026}
\begin{document}
\maketitle
\begin{abstract}
We extend Khesin--Li (2024), Section 8, by proving the operation-generated,
decoration-free graph formulation of Conjecture 8.3 (page 23) and correcting
the shared-output bound in Proposition 8.2 (page 23)~\cite{KL}.
For bipartite encoder graphs in one class, local complementations and
input-edge deletions can be eliminated: output permutations and binary
input-row additions suffice. The output cut-rank function is invariant
under the original operations, up to output relabelling. At a bipartite
endpoint this function determines the binary kernel, hence the row space,
of the input--output matrix. For two inputs, the number of zero columns
and the unordered multiplicities of the three nonzero column types
classify the graphs. Their smallest achievable number of shared outputs
is the least of these multiplicities, at most $\lfloor(n-z)/3\rfloor$
for $z$ zero columns. The $[3,2]$ matrix with all three nonzero column
types has minimum one, whereas the published bound gives zero.
\end{abstract}

\section{The statement and its scope}
Conjecture 8.3 of \cite[Section 8, p.~23]{KL} proposes a reduction
``using only output permutations, input permutations, and row operations''
on the input--output adjacency matrix. Theorem~\ref{thm:main} gives this reduction for the
graph operations specified below. Corollary~\ref{cor:two} then gives a
two-input classification and the exact correction to the adjacent bound.

Let $I=\{1,\ldots,k\}$ and $O=\{1,\ldots,n\}$ be disjoint labelled sets,
where $k,n\geq0$. All graphs are finite simple graphs on $I\sqcup O$.
A \emph{bipartite encoder graph} has edges only between $I$ and $O$;
its input--output matrix is $B_G\in\F^{k\times n}$.
Define $G\sim H$ to be the equivalence relation generated by:
\begin{enumerate}
\item local complementation at any vertex (toggle the edges between its distinct neighbours);
\item arbitrary output permutations and arbitrary input permutations;
\item adding the input--output row of input $j$ to that of input $i$, for $i\ne j$, over $\F$, leaving edges within $I$ and within $O$ unchanged;
\item deletion of an input--input edge.
\end{enumerate}
The equivalence closure allows either direction of each operation.
Intermediate graphs need not be bipartite.
These are the graph operations of \cite[Conjecture 6.3]{KL}, with output
Clifford decorations omitted. The paper uses such operations to construct
the classes in Section 7 and poses the reduction of moves in Conjecture 8.3,
on page 23. The same conjecture appears as Conjecture 5.8.3 in
\cite[Chapter 5]{Khesin}.
The earlier PRIMES version \cite[Section 7, Conjecture 7.1, p.~14]{KL0}
poses the same reduction. Its numbering differs from \cite{KL}, where
Conjecture 7.1 concerns class-size divisibility.
An earlier report by Li, dated 19 August 2023, already poses this
reduction as Conjecture H \cite[Section VII, printed p.~20]{Li}.

\begin{theorem}\label{thm:main}
If $G$ and $H$ are bipartite encoder graphs and $G\sim H$, then a finite
sequence of output permutations, input permutations and additions of
distinct input rows takes $G$ to $H$. Every graph in this sequence is
bipartite.
\end{theorem}

The theorem concerns the specified graph relation. Identifying that relation
with the quantum-code equivalence of \cite[Definition 6.2]{KL} requires
the separate completeness assertion in its Conjecture 6.3.
No rank, connectedness, pivot or normalization condition is imposed here.

\section{Cut ranks retain the binary kernel}
For any graph $G$ and any $S\subseteq O$, put
\[
 \rho_G(S)=\rk_{\F} A_G[S,(I\sqcup O)\setminus S].
\]
Cut rank also measures bipartite graph-state entanglement
\cite[Proposition 3, equation (43)]{Hein}; the argument below uses only
binary matrix operations.
Its invariance under local complementation is classical
\cite[Proposition 2.6, p.~8]{Oum}. We include that argument and check
the additional encoder generators explicitly.

\begin{lemma}\label{lem:invariant}
If $G\sim H$, some output relabelling $\sigma$ of $G$ satisfies
$\rho_{G^\sigma}(S)=\rho_H(S)$ for every $S\subseteq O$.
\end{lemma}
\begin{proof}
Fix $S$ and let $M=A_G[S,(I\sqcup O)\setminus S]$.
If a local complementation is performed at $v\in S$, then for each
$a\in S\setminus\{v\}$ adjacent to $v$, add row $v$ of $M$ to row $a$.
These invertible operations give the new cut matrix: its $(a,b)$ entry
is $M_{ab}+A_G(a,v)A_G(v,b)$, and row $v$ is unchanged.
For $v\notin S$, use the corresponding column additions instead.
Thus local complementation preserves every cut rank.

All inputs lie outside $S$. Adding input row $j$ to input row $i$
therefore adds column $j$ of $M$ to column $i$, an invertible operation.
Deleting an input edge does not affect $M$. Input permutations only
permute its columns. An output permutation transports $S$ and the cut
indices together. Each generator consequently preserves the whole
function up to an output relabelling. The same is true of an inverse,
and composition gives a single relabelling along any finite equivalence
chain.
\end{proof}

If $G$ is bipartite, columns indexed by $O\setminus S$ in its cut matrix
are zero; the other block is $B_G[:,S]^t$. Hence
\begin{equation}\label{eq:profile}
 \rho_G(S)=\rk B_G[:,S].
\end{equation}
For any binary matrix $B$ with $n$ columns, let
\[
 N_B(S)=|\{x\in\ker B:\operatorname{supp}(x)\subseteq S\}|,
 \qquad d_B(S)=\mathbf1_{\{B\mathbf1_S=0\}}.
\]
Here $\mathbf1_S\in\F^n$ is the vector with support exactly $S$.
Rank--nullity and the uniqueness of a binary vector with prescribed
support give
\begin{equation}\label{eq:counts}
 N_B(S)=2^{|S|-\rk B[:,S]}=\sum_{T\subseteq S}d_B(T).
\end{equation}
The first equality follows by extending vectors in $\ker B[:,S]$ by zero.
The second counts each supported binary vector once.

The following is binary-matroid representation uniqueness
\cite[Proposition 6.6.5, p.~209]{Oxley}, attributed there to Brylawski
and Lucas (1976): over $\F$, projective scalings are trivial.
Here is a self-contained supported-vector-counting proof, including
redundant rows.
\begin{lemma}\label{lem:recovery}
For binary matrices $B,C$ of the same size, the following are equivalent:
\[
 \rk B[:,S]=\rk C[:,S]\ (S\subseteq O),\qquad
 \ker B=\ker C,\qquad \Row B=\Row C.
\]
\end{lemma}
\begin{proof}
Equal ranks give equal $N$ by \eqref{eq:counts}. Starting at $S=\varnothing$
and proceeding by cardinality, the identity
\[
 d_B(S)=N_B(S)-\sum_{T\subsetneq S}d_B(T)
\]
recovers every $d_B(S)$ and shows that $d_B=d_C$.
Every vector in $\F^n$ is exactly one $\mathbf1_S$, so the kernels coincide.
Conversely, equal kernels have equal subspaces supported in each $S$,
and \eqref{eq:counts} recovers the ranks. Finally
$\ker B=(\Row B)^\perp$ for the nondegenerate coordinate dot product on
$\F^n$, so equality of kernels is equivalent to equality of row spaces.
\end{proof}

\section{Realizing the restricted moves}
\begin{lemma}\label{lem:realize}
If $B,C\in\F^{k\times n}$ have equal row spaces, there exists
$R\in\GL_k(\F)$ with $C=RB$. Thus row permutations and additions of
distinct rows take $B$ to $C$.
\end{lemma}
\begin{proof}
Regard $B,C$ as maps $\F^n\to\F^k$. Their kernels agree by
Lemma~\ref{lem:recovery}. The map $Bx\mapsto Cx$ is therefore a
well-defined linear isomorphism $\operatorname{im}B\to\operatorname{im}C$:
well-definedness and injectivity both follow from equality of kernels,
and surjectivity holds by definition of the image.
Choose a basis of each image, with corresponding vectors under this
isomorphism, and extend both to bases of $\F^k$. Mapping the extended
bases to each other gives an invertible map $R$ satisfying $RBx=Cx$
for all $x$.

Gaussian elimination reduces an invertible binary square matrix to
the identity using row swaps and additions of a different row.
At a pivot, any nonzero entry is already $1$, so no scaling is needed.
Reversing elimination expresses $R$ as a product of these elementary
matrices. Acting on a bipartite graph, these moves change only its
input--output matrix, so all intermediate graphs remain bipartite.
\end{proof}

\begin{proof}[Proof of Theorem~\ref{thm:main}]
Choose $\sigma$ from Lemma~\ref{lem:invariant}. Equation~\eqref{eq:profile}
gives equal ranks on every column subset for $B_{G^\sigma}$ and $B_H$.
Lemma~\ref{lem:recovery} gives equal row spaces, and
Lemma~\ref{lem:realize} supplies the required row moves after the output
relabeling. Empty input or output sets cause no exception: their
bipartite graphs are edgeless, and the empty sequence suffices.
\end{proof}

\begin{corollary}\label{cor:criterion}
For bipartite $G,H$, the following conditions are equivalent:
\begin{enumerate}
\item $G\sim H$;
\item some output relabelling $\sigma$ gives
$\Row B_{G^\sigma}=\Row B_H$;
\item some output relabelling $\sigma$ gives
$\rho_{G^\sigma}(S)=\rho_H(S)$ for every $S\subseteq O$.
\end{enumerate}
\end{corollary}
\begin{proof}
The preceding lemmas prove $1\Rightarrow3\Rightarrow2$.
Lemma~\ref{lem:realize} realizes condition 2 by allowed moves, proving
$2\Rightarrow1$.
\end{proof}

Enumerating the $n!$ output permutations and comparing reduced row
spaces gives a finite decision procedure. It is not a polynomial-time
bound. In matrix language this is permutation equivalence of binary row
spaces. The subset ranks are the rank function of the associated binary
matroid; matroid-circuit encodings remove basis freedom in modern CSS
code-equivalence algorithms \cite[Section 6.4.1]{BSW}. Our additional graph
step is the invariance in Lemma~\ref{lem:invariant}; no quantum-equivalence
completeness is inferred from that algebraic comparison.

Related move-reduction theorems have different hypotheses.
Kim--Oum \cite[Theorem 5.3, p.~7]{KO} show that two bipartite graphs
with equal cut ranks on \emph{every} vertex subset are pivot-equivalent.
Our input-row operations need not preserve cuts meeting the inputs:
$B=(1,0)^t$ changes to $C=(1,1)^t$ by one addition, changing the
second-input cut rank from zero to one. Thus that all-vertex-cut theorem
does not directly apply to our enlarged relation.
Danielsen--Parker \cite[Theorems 1--2, pp.~5--7]{DP} relate edge-local
complementation to classical-code permutation equivalence; their code
associated with a bipartite block $B$ is $\Row[I_k\mid B]$ of length
$k+n$, rather than $\Row B$ of length $n$ here.

Later structural results also address distinct relations.
Wang--Krastanov \cite[Theorem 1 and Corollary 5]{WK} give polynomial-time
local-Clifford conversion to CSS form and recognition of a bipartite
graph in an ordinary local-complementation orbit, without a prescribed
encoder partition. Albert \cite[Appendix F.2, Lemma F.2]{Albert} reduces
local symplectic maps between CSS-form subspaces to subset Hadamards.
Neither statement includes the additional fixed-partition input-row
and input-edge generators used in Theorem~\ref{thm:main}.

\begin{corollary}
Input permutations can be omitted from the restricted generating set.
\end{corollary}
\begin{proof}
the three additions $(r_i,r_j)\mapsto(r_i+r_j,r_j)
\mapsto(r_i+r_j,r_i)\mapsto(r_j,r_i)$ exchange two rows.
\end{proof}

\section{Two inputs: classification and the sharp bound}
For a binary $2\times n$ matrix, write $z$ for the number of zero columns
and $a,b,c$ for the multiplicities of $(1,0)^t,(0,1)^t,(1,1)^t$.
A shared output is a column of the last type. In a pivot form the two
pivot columns have types $(1,0)^t$ and $(0,1)^t$, so all shared outputs
are free outputs in the terminology of \cite[Section 8]{KL}.

\begin{corollary}\label{cor:two}
Two bipartite graphs with two inputs and $n$ outputs are equivalent if
and only if they have the same $z$ and the same unordered triple
$\{a,b,c\}$, with multiplicities retained. Over all equivalent bipartite
forms, the minimum number of shared outputs is
\[
 m=\min(a,b,c)\leq\left\lfloor\frac{n-z}{3}\right\rfloor.
\]
The minimum total number of edges is $n-z+m$.
When the row rank is two, the minimum can be attained in a pivot form.
The upper bound is sharp for every $n-z\geq2$.
For two inputs, output cut ranks on singletons and pairs already determine
the row space, with the output labels fixed.
\end{corollary}
\begin{proof}
The action of $\GL_2(\F)$ on the three nonzero vectors is the full
symmetric group: an ordered pair of distinct nonzero vectors is a basis,
and specifying its image determines an invertible map. There are six such
maps. They fix zero and permute $a,b,c$ arbitrarily. Thus the row-space
criterion in Corollary~\ref{cor:criterion}, followed by a column permutation,
is precisely equality of $z$ and the unordered triple. It also makes the
possible shared-output counts exactly $a,b,c$. Their minimum is at most
their average. A nonzero column contributes one edge, and a shared column
one additional edge, giving the edge minimum. If the rank is two, at least two types occur; after assigning
the least frequent type to $(1,1)^t$, both other types still occur and supply
the two pivot columns. For sharpness, distribute $n-z$ columns among the
three types with counts differing by at most one. For $n-z\geq2$ the
resulting matrix has rank two and minimum $\lfloor(n-z)/3\rfloor$.
Finally, singleton ranks detect zero columns; for nonzero columns, a pair
has rank one exactly when the two columns coincide. These ranks recover
the partition into the at most three nonzero types, and any relabelling
of those types is realized by $\GL_2(\F)$. Thus they recover the row space.
\end{proof}

This reproves the three-choice observation already in Li's report
\cite[Section VII, p.~20, after Claim G]{Li}, and later following
\cite[Proposition 8.2, p.~23]{KL}: for independent rows $u,v$, the
unordered bases are $\{u,v\}$, $\{u,u+v\}$ and $\{v,u+v\}$.
Corollary~\ref{cor:two} turns that observation into a classification:
counting the four column types and sorting three integers takes $O(n)$
operations. It answers the preceding equivalence-testing question for two
inputs, without asserting a polynomial-time algorithm for variable $k$.
It also selects a canonical representative: if $a\leq b\leq c$ are the
sorted nonzero-type counts, use $b$ copies of $(1,0)^t$, $c$ copies of
$(0,1)^t$, $a$ copies of $(1,1)^t$, and $z$ zero columns, in that order.
For full rank, $b,c>0$ and one column of each of the first two types can
instead be placed first as the pivots.
This representative also minimizes the edge count, addressing the choice
made after \cite[Proposition 8.1, p.~22]{KL}.

\begin{corollary}
For one input, the row weight classifies bipartite graphs under $\sim$.
\end{corollary}
\begin{proof}
The only invertible binary one-by-one matrix is the identity. By
Corollary~\ref{cor:criterion}, equivalence is therefore output-permutation
equivalence of the rows, which is exactly equality of their weights.
\end{proof}

\begin{corollary}\label{cor:connected}
A full-rank two-input bipartite graph has no disconnected graph in its
operation-generated equivalence class exactly when $z=0$ and $a,b,c>0$.
Its classes with this property correspond to unordered triples of positive
integers summing to $n$. In particular there is exactly one such class
for $n=4$, already in \cite[Claim F, pp.~17--18]{Li} and later
\cite[Proposition 8.1, p.~22]{KL}.
\end{corollary}
\begin{proof}
If $z>0$, an output is isolated. If a nonzero type is absent, an invertible
row change sends the two occurring types to $(1,0)^t,(0,1)^t$ and disconnects
the graph into two stars. Conversely, suppose an equivalent graph $H$ is
disconnected. Its components partition the inputs into sets $I_j$ and
the outputs into sets $S_j$. Output cut-rank invariance, after relabelling,
gives
\[
 \rk B[:,S_j]=\rho_H(S_j)\leq |I_j|,
 \qquad \rk B=\rho_H(O)=2.
\]
Since the $S_j$ partition the columns and $\sum_j|I_j|=2$,
subadditivity of rank forces equality throughout
$2\leq\sum_j\rk B[:,S_j]\leq2$.
A component with no input would therefore give zero columns, contrary
to $z=0$; a component with inputs but no outputs is impossible by the
same equality. Thus disconnectedness requires two components, each with
one input, and each column block has rank one. Each block has at most one
nonzero column type, contradicting the presence of all three types.
The classification now follows from Corollary~\ref{cor:two}; for $n=4$
the sole positive triple is $\{1,1,2\}$.
\end{proof}

Outside the normalized setting, connectedness of an individual drawing
is not invariant: the disconnected graph with matrix $I_2$ becomes
connected after adding its second row to the first,
$\bigl(\begin{smallmatrix}1&1\\0&1\end{smallmatrix}\bigr)$ being a four-vertex
path. The latter matrix is not in RREF. This unrestricted assertion is
not made in \cite{KL0}, whose Section 5 fixes RREF before its discussion
of connectivity. Corollary~\ref{cor:connected} states the invariant
class condition without an RREF hypothesis.

\paragraph{Correction to Proposition 8.2.}
The bound in \cite[Proposition 8.2, p.~23]{KL} is
$\lfloor n/2-1\rfloor$. For
\[
 B=\begin{pmatrix}1&0&1\\0&1&1\end{pmatrix},
 \qquad n=3,
\]
the three multiplicities are all one. Every equivalent bipartite form
therefore has one shared output, not at most zero. This is a full-rank,
pivot-form example satisfying the encoder hypotheses; it is also connected
throughout its graph-operation class by Corollary~\ref{cor:connected}.
Corollary~\ref{cor:two} supplies the corrected sharp bound for all $n\geq2$.
The published numerical bound does hold for $n\geq4$, because
$\lfloor n/3\rfloor\leq\lfloor n/2\rfloor-1$ there. In the displayed
inequality in its proof, $n-1-\lfloor n/2\rfloor<\lfloor n/2\rfloor$
is not strict for odd $n$; the exact minimum above avoids that step.
The same correction applies to \cite[Claim 7.2, p.~14]{KL0} and
\cite[Proposition 5.8.2, pp.~127--128]{Khesin}, and to
\cite[Claim G, pp.~18--20]{Li}, which state the same bound.

The same example also obstructs a zero-shared-output form under product
output unitaries and permutations, independently of Conjecture 6.3.
For a bipartite graph, its encoding matrix has entries proportional to
$(-1)^{u^tBx}$, so a Hadamard on each output identifies its codespace with
$\operatorname{span}\{|y\rangle:y\in\Row B\}$.
For the displayed $B$ this is the even-parity code, with projector
$P=(1+Z\otimes Z\otimes Z)/2$. A full-rank $2\times3$ matrix with no
shared output has either a zero column or two equal columns; its
one-dimensional binary kernel is generated by a vector of weight one or
two. Its projector therefore has the form $(1+Z^w)/2$, $1\leq|w|\leq2$.
Every one-qubit partial trace of $2P-1=Z\otimes Z\otimes Z$ vanishes,
whereas tracing out a qubit outside the support of $Z^w$ is nonzero.
Product unitaries and permutations preserve this property. Input
unitaries do not change the codespace. Hence these projectors cannot be
equivalent under those transformations.

\section{Boundaries and coverage}
The following extensions are not asserted in the source. Dropping the
bipartite endpoint hypothesis fails: with two inputs and no outputs, the
single-edge graph and the edgeless graph are related by input-edge
deletion. The restricted moves cannot connect them: input permutations
preserve the edge, and input-row additions have no output entries to change.

Unchanged kernel recovery over other fields also fails. Over $\mathbb F_3$, $(1,1)$ and
$(1,2)$ have identical ranks on all column subsets, but $(1,2)$ as a
column vector is in the kernel of the first and not of the second.
Even over $\F$, the total rank alone is insufficient: $(1,0)$ and $(1,1)$
both have rank one, but only the first has a zero column.

The entanglement suggestion after \cite[Conjecture 8.3, pp.~23--24]{KL}
is realized here by cut ranks. For two inputs, singleton and pair output
cuts suffice by Corollary~\ref{cor:two}. Extending that completeness to
three inputs is not asserted in the source and fails: the matrices with
columns $(e_1,e_2,e_3,e_1+e_2)$ and $(e_1,e_2,e_3,e_1+e_2+e_3)$ have
singleton ranks one and pair ranks two, but their one-dimensional kernels
have nonzero vectors of weights three and four, respectively. No output
permutation equates those kernels, so the graphs are inequivalent.

\begin{center}
\small
\begin{tabular}{@{}p{0.30\linewidth}p{0.64\linewidth}@{}}
\toprule
Source item & Coverage\\
\midrule
\cite[Conjecture 8.3]{KL}; \cite[Conjecture 5.8.3]{Khesin};
\cite[Conjecture 7.1]{KL0}; \cite[Conjecture H]{Li}
 & Proved for the operation-generated, decoration-free graph relation
   (Theorem~\ref{thm:main}).\\[3pt]
Equivalence-testing proposal before \cite[Conjecture 8.3, p.~23]{KL}
 & Answered: yes for two inputs, with the linear-time criterion above;
   general-$k$ finite criterion in Corollary~\ref{cor:criterion}, without
   a polynomial-time bound.\\[3pt]
Three input bases after \cite[Proposition 8.2]{KL}
 & Reproved; earlier statement Li (2023), after Claim G; later
   Khesin--Li (2024), also after Claim 7.2 in \cite{KL0}.\\[3pt]
\cite[Proposition 8.1, p.~22]{KL}
 & Reproved; prior proof Li (2023), Claim F; later Khesin--Li (2024), also \cite[Claim 7.1]{KL0}. Corollary~\ref{cor:connected}
   gives the classification for every $n$.\\[3pt]
\cite[Proposition 8.2, p.~23]{KL}; \cite[Proposition 5.8.2]{Khesin};
\cite[Claim 7.2]{KL0}; \cite[Claim G]{Li}
 & Refuted-and-corrected: $n=3$ is a counterexample; the exact minimum
   and sharp replacement are in Corollary~\ref{cor:two}.\\[3pt]
Unrestricted drawing-connectedness invariance
 & Not asserted in the source's normalized setting; fails for $I_2$
   after one row addition. The class criterion is Corollary~\ref{cor:connected}.\\[3pt]
\cite[Conjectures 6.3 and 7.1]{KL}
 & Not covered: quantum-equivalence completeness and normalized-class
   divisibility, respectively.\\[3pt]
Small-$n$ class-count patterns, \cite[Section 7]{KL}
 & Not covered; counting bipartite-containing classes does not count
   all classes or their normalized members.\\[3pt]
Entanglement proposal after \cite[Conjecture 8.3]{KL}
 & All-size proof by output cut ranks; singleton and pair cuts suffice
   for two inputs, not for three.\\[3pt]
Canonical-choice discussion, \cite[Section 8, pp.~21--22]{KL}
 & Two-input representative constructed above; general canonical
   efficiency not covered.\\[3pt]
Classes without bipartite forms, \cite[Section 8, pp.~23--24]{KL}
 & Not covered; the theorem assumes bipartite endpoints.\\[3pt]
Extensions omitting bipartiteness or replacing $\F$ by $\mathbb F_3$
 & Not asserted in the source; fail for the examples above.\\
\bottomrule
\end{tabular}
\end{center}

\paragraph{Verification.}
The \href{https://github.com/the-omega-institute/trureturing/blob/5a1fca9e717166d9d52651194f897dda2b9ff37f/D5/S3/Quantum/Codes/KhesinLiBipartiteEncoderEquivalence.lean}{pinned bipartite-encoder module}
contains declaration \texttt{result : claim}, proving exactly the implication
in Theorem~\ref{thm:main} for every natural $k,n$.
Its \href{https://github.com/the-omega-institute/trureturing/blob/5a1fca9e717166d9d52651194f897dda2b9ff37f/D5/S3/Quantum/Codes/KhesinLiEquivalenceClasses.lean}{imported \texttt{Step} definition}
is the operation relation specified above; it also defines
\texttt{addInputRow}. The equivalence criterion, two-input classification,
sharp bound, connectedness criterion, correction and boundary examples
have written proofs here; they are not conclusions of \texttt{result : claim}.
The accompanying \texttt{verify.py} checks finite binary matrices and
complete graph-operation components in the ranges in \texttt{verification.json}.
It also enumerates all $21\,845$ binary $2\times n$ matrices with
$0\leq n\leq7$, all six input-basis changes for each matrix, and the
three-output counterexample. These finite checks are distinct from the
all-size written proofs.

\begin{thebibliography}{99}
\bibitem{KL} A. B. Khesin and A. M. Li,
\emph{Equivalence Classes of Quantum Error-Correcting Codes},
arXiv:2406.12083v1 (2024), Sections 6--8.
\href{https://arxiv.org/abs/2406.12083v1}{arXiv:2406.12083v1}.
\bibitem{Khesin} A. B. Khesin,
\emph{Quantum Computing from Graphs},
Ph.D. thesis, Massachusetts Institute of Technology (2025), Chapter 5;
\href{https://arxiv.org/abs/2501.17959v1}{arXiv:2501.17959v1}.
\bibitem{Hein} M. Hein, J. Eisert and H. J. Briegel,
\emph{Multi-party entanglement in graph states},
Physical Review A \textbf{69}, 062311 (2004);
\href{https://arxiv.org/abs/quant-ph/0307130v7}{arXiv:quant-ph/0307130v7},
Proposition 3 and equation (43).
\bibitem{BSW} M. Bichmann, L. Schmid and R. Wille,
\emph{Methods for Automated Equivalence Checking of Quantum
Error-Correction Codes},
\href{https://arxiv.org/abs/2610.18218v1}{arXiv:2610.18218v1} (2026),
Sections 3--4 and 6.4.1.
\bibitem{KL0} A. B. Khesin and A. M. Li,
\emph{Canonical Forms and Equivalence Classes of QECC's in ZX Calculus},
MIT PRIMES (4 February 2024), Sections 7--8, pp.~13--15.
\href{https://math.mit.edu/research/highschool/primes/materials/2023/Khesin-Li.pdf}{PRIMES paper};
dated entry 391 in the
\href{https://math.mit.edu/research/highschool/primes/papers.html}{PRIMES papers list}.
\bibitem{Li} A. M. Li,
\emph{Canonical Forms and Equivalence Classes of Quantum Clifford Codes
in ZX Calculus}, report dated 19 August 2023, Section VII, pp.~17--20.
\href{https://www.yau-science-awards.org/upload/202311/13/202311131114411646/202311131114411646.pdf}{Yau report}.
\bibitem{Oum} S. Oum, \emph{Rank-width and vertex-minors},
Journal of Combinatorial Theory, Series B \textbf{95} (2005), 79--100,
Proposition 2.6 (author version p.~8).
\href{https://mathsci.kaist.ac.kr/~sangil/pdf/vertexminor-revision.pdf}{Author version}.
\bibitem{Oxley} J. Oxley, \emph{Matroid Theory}, second edition,
Oxford University Press (2011), Proposition 6.6.5, p.~209.
\href{https://www.math.ens.psl.eu/~benoist/refs/Oxley.pdf}{Text}.
\bibitem{KO} D. Kim and S. Oum, \emph{Vertex-minors of graphs: A survey},
Discrete Applied Mathematics \textbf{351} (2024), 54--73;
author version dated 18 March 2024, Theorem 5.3, p.~7.
\href{https://sangil.dimag.kr/legacy/2024/03/2023vertexminors-survey-revised.pdf}{Author version}.
\bibitem{DP} L. E. Danielsen and M. G. Parker,
\emph{Edge Local Complementation and Equivalence of Binary Linear Codes},
Designs, Codes and Cryptography \textbf{49} (2008), 161--170;
\href{https://arxiv.org/abs/0710.2243v4}{arXiv:0710.2243v4}, Theorems 1--2.
\bibitem{WK} M. Wang and S. Krastanov,
\emph{Local Clifford equivalence to CSS codes: Polynomial-time classification
and conversion}, \href{https://arxiv.org/abs/2610.28737v1}{arXiv:2610.28737v1}
(2026), Theorem 1 and Corollary 5.
\bibitem{Albert} V. V. Albert,
\emph{Beyond transversality: structure of Clifford circuits for CSS codes},
\href{https://arxiv.org/abs/2608.05688v2}{arXiv:2608.05688v2} (2026),
Appendix F.2, Lemma F.2.
\end{thebibliography}
\smallskip
\textcopyright\ 2026 Haobo Ma and Wenlin Zhang.
Licensed under \href{https://creativecommons.org/licenses/by/4.0/}{Creative
Commons Attribution 4.0 International (CC BY 4.0)}.
\end{document}
